\section{二八〇年的统一与它尚未完成的工作}
\label{sec:western-jin-unification-fracture}

\begin{event}{\eventanchor{TK-0280-SUN-HAO-SURRENDER}{孙皓降服与西晋统一}}{太康元年三月（二八〇年）}{西晋与孙吴 · 建业}{降服可核，战役规模与接收范围不定}
《晋书·武帝纪》《晋书·王濬传》《三国志·孙皓传》与《建康实录》可互校建业降附和孙吴灭亡的结果。舰只、兵员、各路抵达先后和战后安置的范围在传世叙事中并不一致，不能由统一的名义倒推江南已经完成整合。
\end{event}

太康元年三月，长江下游的战事走到建业城下。益州刺史王濬所部顺流而下，荆州与江淮诸军从别路施压；孙皓出降，\hyperlink{evt:TK-0280-SUN-HAO-SURRENDER}{西晋·孙吴降服}。这一节点可以由《晋书·武帝纪》《王濬传》、《三国志·孙皓传》与《建康实录》相互校读：孙吴灭亡、建业降附和西晋完成最高政权的统一，已有可靠的年序和结果。

不过，沿江推进的兵额、舰只数目、各路抵达的先后，以及王濬、王浑各自应得多少战功，在传世叙事中并不一致。它们不能被画成一张精确的作战图，更不能让胜者传记替代江东每一处城邑的经验。可把握的是，建业一降，吴的州郡、官民和水陆节点必须改向洛阳的任命与征发体系；如何安置降军、沿用旧吏、核对簿籍，史书留下的却远少于降表与将领姓名。

这一区分改变了“统一”的读法。舰队能够通过长江，说明晋军取得了水道和军事上的优势；一县的税粮是否照旧入仓、一支旧军是否接受新将、渡口由谁巡护，则是此后才会逐处发生的接收工作。没有材料把这些过程完整记录下来，不能把二八〇年写成江南立刻与洛阳采用同一套日常秩序。统一首先改变了最高裁决和官号的归属，随后才打开了重新安排人、文书和供给的漫长事务。

\begin{causalchain}[建业降服所留下的工作]
\textbf{直接结果}：三国并立的最高统治权结束，西晋取得吴的州郡与人口。\\
\textbf{使接收成为可能的条件}：沿江军事推进使建业在政治和军事上孤立。\\
\textbf{不能由降服推出的结论}：江南的赋役、地方关系与社会秩序已经无缝改造。\\
\textbf{后续压力}：新朝必须让任命、簿籍、守备和水路供给在不同地区持续运转。
\end{causalchain}

次年，离建业很远的汲郡发生了另一件进入洛阳的事。墓中竹书被发现后，朝廷命荀勖、和峤等参与整理，\eventanchor{JH-0281-JI-TOMB-BAMBOO}{西晋·汲冢竹书}。出土和入藏是可确认的；墓主的认定、竹书原有的篇章形态，以及今本《竹书纪年》同原简的关系，仍不能强作结论。这里没有必要把一批竹简附会成统一的“文化象征”：它同建业降服之间并无可证的直接因果。

它所显示的是另一种接收。竹简从地方墓葬进入官府，须经识读古字、分类、校勘与编目，才会成为朝廷能够引用的知识。西晋的统一不只涉及谁在江面上获胜，也涉及洛阳怎样把来自不同地方、不同年代的材料纳入自己的档案与判断。可是，整理官署的活动不能证明整个帝国已有同样的文化或行政能力；现存材料只够让我们看见这一批文字的入藏和参与者，不能补出未留下记录的地方反应。

从建业的降表到汲郡的竹简，二八〇至二八一年留下的是两条并行而尺度不同的线索：一条是军事胜利后对州郡、人户与水道的接收，另一条是国家将偶然发现的文字变成可用知识的过程。二者都要求中介者、文书和持续的组织，却不能彼此解释。这样的限度也使后来的继承问题更清楚：皇帝的名义可以统一两条线，但未必能保证宫城、官署与掌兵者继续接受同一套裁决。

太熙元年，武帝病重，以外戚杨骏受遗诏辅政，\eventref{JH-0289-YANG-JUN-REGENT}{西晋·杨骏辅政}。受遗辅政和杨骏居要职可由《晋书·武帝纪》《杨骏传》及《资治通鉴·晋纪》确认；遗诏的措辞、武帝意向和其他顾命者的作用则有争议。近十年前的统一并没有替这一时刻预先写好答案：继承如何把诏令、尚书官署和宫城卫队维系在一起，成为下一节的核心。

\begin{sourcenote}[建业与汲郡各能证实到哪里]
二八〇年的年序与政权转换主要依《晋书》《三国志》《建康实录》交叉校读；它们可确认降服和统一的结果，不足以精算兵数、路线、安置范围或人物动机。汲冢竹书的发现与入藏可由《晋书·武帝纪》《束皙传》《荀勖传》追踪，但传本、墓主与篇章关系仍须由出土文献和战国史书重构研究继续检验。两组材料都不能替代江东或汲郡社会的全景。
\end{sourcenote}


% 年序补线：本节将既有账本中尚未初次落入正文的条目接入相应历史场景。
\subsection{统一之后：水道、官署与继承的接合}
对于太康二年（281年）的西晋而言，\eventanchor{JH-LEDGER-001-1}{汲郡魏襄王墓出土竹书，朝廷命荀勖、和峤等整理其书写材料}、\eventanchor{JH-LEDGER-001-2}{汲郡魏襄王墓出土竹书，朝廷命荀勖、和峤等整理其书写材料}、\eventanchor{JH-LEDGER-001-3}{汲郡魏襄王墓出土竹书，朝廷命荀勖、和峤等整理其书写材料}、\eventanchor{JH-LEDGER-001-4}{汲郡魏襄王墓出土竹书，朝廷命荀勖、和峤等整理其书写材料}、\eventanchor{JH-LEDGER-001-5}{汲郡魏襄王墓出土竹书，朝廷命荀勖、和峤等整理其书写材料}、\eventanchor{JH-LEDGER-001-6}{“汲郡魏襄王墓出土竹书，朝廷命荀勖、和峤等整理其书写材料”的异源材料与影响范围的复核}都回到同一具体节点：汲郡魏襄王墓出土竹书，朝廷命荀勖、和峤等整理其书写材料。 现有材料只能把这一变化放在连续年序中；前期条件、命令响应、地方执行、后续调整和异文问题是同一事件的不同读法，不能伪装成六场独立历史。 《晋书·武帝纪》；《晋书·束皙传》；《荀勖传》；汲郡魏襄王墓出土竹书提供该簇的最低证据链；核心事件已有传世材料可供核对；各个观察面向的参与者、执行范围和时间先后仍须逐案核验。

太熙元年（289年）的西晋，通过\eventanchor{JH-LEDGER-002-1}{武帝病重时以外戚杨骏受遗诏辅政，围绕太子司马衷的继承安排集中于宫廷}、\eventanchor{JH-LEDGER-002-2}{武帝病重时以外戚杨骏受遗诏辅政，围绕太子司马衷的继承安排集中于宫廷}、\eventanchor{JH-LEDGER-002-3}{武帝病重时以外戚杨骏受遗诏辅政，围绕太子司马衷的继承安排集中于宫廷}、\eventanchor{JH-LEDGER-002-4}{武帝病重时以外戚杨骏受遗诏辅政，围绕太子司马衷的继承安排集中于宫廷}、\eventanchor{JH-LEDGER-002-5}{武帝病重时以外戚杨骏受遗诏辅政，围绕太子司马衷的继承安排集中于宫廷}、\eventanchor{JH-LEDGER-002-6}{“武帝病重时以外戚杨骏受遗诏辅政，围绕太子司马衷的继承安排集中于宫廷”的异源材料与影响范围的复核}把武帝病重时以外戚杨骏受遗诏辅政，围绕太子司马衷的继承安排集中于宫廷这一历史过程从条件、处置到余波连成一体。 现有材料只能把这一变化放在连续年序中；前期条件、命令响应、地方执行、后续调整和异文问题是同一事件的不同读法，不能伪装成六场独立历史。 《晋书·武帝纪》；《晋书·杨骏传》；《资治通鉴·晋纪》提供该簇的最低证据链；核心事件已有传世材料可供核对；各个观察面向的参与者、执行范围和时间先后仍须逐案核验。



\subsection{280—289年：材料所见的连续压力}
统一后江东州郡的接收仍在调整；诏令须穿过旧水路与地方吏治，户籍、仓储和交通只能逐次对接。

沿着公元280年（年序桥接），西晋与江东州郡的\eventanchor{JH-DENS-280-A}{西晋接收江东后，州郡官员、户籍和漕运安排仍在调整}指向西晋接收江东后，州郡官员、户籍和漕运安排仍在调整，\eventanchor{JH-DENS-280-B}{西晋与江东州郡的朝廷和州郡围绕户籍、田租、军粮与转运的安排进入材料}则把西晋与江东州郡的朝廷和州郡围绕户籍、田租、军粮与转运的安排进入材料接入同一条年序。 这使读者能把下一年的军事、行政或社会变化放回尚未消散的限制中。 就可说范围而言，前一判断止于《晋书·武帝纪》《晋书·地理志》；《资治通鉴·晋纪》，本纪与编年材料主要确证政权年序和精英行动；对基层执行、疆界和人口数量的沉默不得以叙事补足；后一判断止于《晋书·食货志》；墓志与区域考古，志书、诏令与本纪可显示制度设计或战争需求，不能直接换算赋额、仓储或实际负担。

在公元281年（年序桥接）的并行行动者之间，\eventanchor{JH-DENS-281-A}{西晋接收江东后，州郡官员、户籍和漕运安排仍在调整}将西晋接收江东后，州郡官员、户籍和漕运安排仍在调整落到时间位置，\eventanchor{JH-DENS-281-B}{西晋与江东州郡的流民、侨置户、军镇居民或地方家族的安置与编户问题构成社会背景}从西晋与江东州郡的流民、侨置户、军镇居民或地方家族的安置与编户问题构成社会背景补足资源与空间的视角。 如此处理不是用空白填满史实，而是避免后来的转折抹去本年仍在运作的条件。 材料边界须分别保留：《晋书·武帝纪》《晋书·地理志》；《资治通鉴·晋纪》与《晋书·食货志》；墓志与区域考古所见不同，且本纪与编年材料主要确证政权年序和精英行动；对基层执行、疆界和人口数量的沉默不得以叙事补足；墓志、家族文本和侨置名目各自有选择性，不能据此统计迁徙规模或概括整体社会态度。

公元282年（年序桥接），\eventanchor{JH-DENS-282-A}{西晋接收江东后，州郡官员、户籍和漕运安排仍在调整}和\eventanchor{JH-DENS-282-B}{西晋与江东州郡的寺院、僧侣、道士和施主网络在材料中连接都城、州郡与边地}共同避免把西晋与江东州郡压缩为单一都城事件：前者关注西晋接收江东后，州郡官员、户籍和漕运安排仍在调整，后者关注西晋与江东州郡的寺院、僧侣、道士和施主网络在材料中连接都城、州郡与边地。 年序因此不把大事件当作一切的终点，而让制度、交通和社会关系继续承受压力。 此处的证据支点是《晋书·武帝纪》《晋书·地理志》；《资治通鉴·晋纪》和《晋书·食货志》；墓志与区域考古；本纪与编年材料主要确证政权年序和精英行动；对基层执行、疆界和人口数量的沉默不得以叙事补足。，同时僧传、道教文本、造像记与遗址的成书或营造年代须分开，个案不能量化信仰或政策落实。

把公元283年（年序桥接）放回连续叙事，\eventanchor{JH-DENS-283-A}{西晋接收江东后，州郡官员、户籍和漕运安排仍在调整}保留西晋接收江东后，州郡官员、户籍和漕运安排仍在调整，\eventanchor{JH-DENS-283-B}{西晋与江东州郡的关隘、渡口、军镇和水陆道路的可通行性继续约束使节、军队和物资的行动}让西晋与江东州郡的关隘、渡口、军镇和水陆道路的可通行性继续约束使节、军队和物资的行动继续约束西晋与江东州郡的行动。 这使读者能把下一年的军事、行政或社会变化放回尚未消散的限制中。 材料边界须分别保留：《晋书·武帝纪》《晋书·地理志》；《资治通鉴·晋纪》与《晋书·食货志》；墓志与区域考古所见不同，且本纪与编年材料主要确证政权年序和精英行动；对基层执行、疆界和人口数量的沉默不得以叙事补足；纪传中的行军路线、兵数与控制范围常被简化或夸饰，地理材料只能提供有限交叉。

\eventanchor{JH-DENS-284-A}{西晋接收江东后，州郡官员、户籍和漕运安排仍在调整}在公元284年（年序桥接）提示西晋接收江东后，州郡官员、户籍和漕运安排仍在调整仍未结案；\eventanchor{JH-DENS-284-B}{西晋与江东州郡的朝廷和州郡围绕户籍、田租、军粮与转运的安排进入材料}以西晋与江东州郡的朝廷和州郡围绕户籍、田租、军粮与转运的安排进入材料说明西晋与江东州郡的选择还受具体条件限制。 如此处理不是用空白填满史实，而是避免后来的转折抹去本年仍在运作的条件。 材料边界须分别保留：《晋书·武帝纪》《晋书·地理志》；《资治通鉴·晋纪》与《晋书·食货志》；墓志与区域考古所见不同，且本纪与编年材料主要确证政权年序和精英行动；对基层执行、疆界和人口数量的沉默不得以叙事补足；志书、诏令与本纪可显示制度设计或战争需求，不能直接换算赋额、仓储或实际负担。

公元285年（年序桥接），西晋与江东州郡的\eventanchor{JH-DENS-285-A}{西晋接收江东后，州郡官员、户籍和漕运安排仍在调整}把西晋接收江东后，州郡官员、户籍和漕运安排仍在调整留在年序中；\eventanchor{JH-DENS-285-B}{西晋与江东州郡的流民、侨置户、军镇居民或地方家族的安置与编户问题构成社会背景}则从西晋与江东州郡的流民、侨置户、军镇居民或地方家族的安置与编户问题构成社会背景观察同一年的约束。 这使读者能把下一年的军事、行政或社会变化放回尚未消散的限制中。 证据不许把这两个观察写成全域事实：《晋书·武帝纪》《晋书·地理志》；《资治通鉴·晋纪》受本纪与编年材料主要确证政权年序和精英行动；对基层执行、疆界和人口数量的沉默不得以叙事补足。限制，《晋书·食货志》；墓志与区域考古亦受墓志、家族文本和侨置名目各自有选择性，不能据此统计迁徙规模或概括整体社会态度限制。

在公元286年（年序桥接），\eventanchor{JH-DENS-286-A}{西晋接收江东后，州郡官员、户籍和漕运安排仍在调整}所见的西晋接收江东后，州郡官员、户籍和漕运安排仍在调整与\eventanchor{JH-DENS-286-B}{西晋与江东州郡的寺院、僧侣、道士和施主网络在材料中连接都城、州郡与边地}所见的西晋与江东州郡的寺院、僧侣、道士和施主网络在材料中连接都城、州郡与边地并行发生在西晋与江东州郡的历史空间。 这两项观察不宣称另有一场未载战争或政策，只把已有压力的时间位置和可见面向分开。 就可说范围而言，前一判断止于《晋书·武帝纪》《晋书·地理志》；《资治通鉴·晋纪》，本纪与编年材料主要确证政权年序和精英行动；对基层执行、疆界和人口数量的沉默不得以叙事补足；后一判断止于《晋书·食货志》；墓志与区域考古，僧传、道教文本、造像记与遗址的成书或营造年代须分开，个案不能量化信仰或政策落实。

\eventanchor{JH-DENS-287-A}{西晋接收江东后，州郡官员、户籍和漕运安排仍在调整}和\eventanchor{JH-DENS-287-B}{西晋与江东州郡的关隘、渡口、军镇和水陆道路的可通行性继续约束使节、军队和物资的行动}使公元287年（年序桥接）的西晋与江东州郡既保留西晋接收江东后，州郡官员、户籍和漕运安排仍在调整，也不略去西晋与江东州郡的关隘、渡口、军镇和水陆道路的可通行性继续约束使节、军队和物资的行动。 年序因此不把大事件当作一切的终点，而让制度、交通和社会关系继续承受压力。 材料边界须分别保留：《晋书·武帝纪》《晋书·地理志》；《资治通鉴·晋纪》与《晋书·食货志》；墓志与区域考古所见不同，且本纪与编年材料主要确证政权年序和精英行动；对基层执行、疆界和人口数量的沉默不得以叙事补足；纪传中的行军路线、兵数与控制范围常被简化或夸饰，地理材料只能提供有限交叉。

沿着公元288年（年序桥接），西晋与江东州郡的\eventanchor{JH-DENS-288-A}{西晋接收江东后，州郡官员、户籍和漕运安排仍在调整}指向西晋接收江东后，州郡官员、户籍和漕运安排仍在调整，\eventanchor{JH-DENS-288-B}{西晋与江东州郡的朝廷和州郡围绕户籍、田租、军粮与转运的安排进入材料}则把西晋与江东州郡的朝廷和州郡围绕户籍、田租、军粮与转运的安排进入材料接入同一条年序。 它们不是由年份自动推出的因果，而是让前后事件之间的资源、道路、户籍或地方网络保持可见。 前者依据《晋书·武帝纪》《晋书·地理志》；《资治通鉴·晋纪》；本纪与编年材料主要确证政权年序和精英行动；对基层执行、疆界和人口数量的沉默不得以叙事补足。后者参照《晋书·食货志》；墓志与区域考古；志书、诏令与本纪可显示制度设计或战争需求，不能直接换算赋额、仓储或实际负担。

在公元289年（年序桥接）的并行行动者之间，\eventanchor{JH-DENS-289-A}{西晋接收江东后，州郡官员、户籍和漕运安排仍在调整}将西晋接收江东后，州郡官员、户籍和漕运安排仍在调整落到时间位置，\eventanchor{JH-DENS-289-B}{西晋与江东州郡的流民、侨置户、军镇居民或地方家族的安置与编户问题构成社会背景}从西晋与江东州郡的流民、侨置户、军镇居民或地方家族的安置与编户问题构成社会背景补足资源与空间的视角。 它们将政治时间同资源和地方网络并置，却不把并置误作已被证实的直接因果。 前者依据《晋书·武帝纪》《晋书·地理志》；《资治通鉴·晋纪》；本纪与编年材料主要确证政权年序和精英行动；对基层执行、疆界和人口数量的沉默不得以叙事补足。后者参照《晋书·食货志》；墓志与区域考古；墓志、家族文本和侨置名目各自有选择性，不能据此统计迁徙规模或概括整体社会态度。


《晋书·武帝纪》、志书与《资治通鉴·晋纪》提供年序，不能量出各地征收、迁徙或寺院网络。
